% Section content template for individual Educative lessons
% This file contains the actual content for a specific section
% Generated from Educative API section content

% Section: Summary: Designing the Amazon Locker Service
% Section ID: 6564313689882624
% Section Slug: summary-designing-the-amazon-locker-service
% Generated: 2025-12-11T15:17:29.890032

Now that you've completed the Amazon Locker service case study, let's take a moment to reflect on and consolidate what we've learned. We 'll revisit the key system requirements, identify the core classes along with their responsibilities and relationships, and highlight the major design principles applied. We
'll also examine how objects interact within the system and walk through the overall workflow to understand how the components come together to achieve the desired functionality.

\subsection{Key requirements}\label{yZtfiqtKpJEGLLknIWleC}
\\
Here is a list of the primary functional requirements for the Amazon Locker system.

\begin{enumerate}
\item
\protect\phantomsection\label{wHeW4hyv9DUPYiKW7Pv2G}
Customers can select a preferred locker location for order pickup during checkout.
\item
\protect\phantomsection\label{6ShAWi6PVxmzin1fiZ1tW}
Orders can contain multiple items packaged together based on locker size availability.
\item
\protect\phantomsection\label{stsiJaIKcA3DXgH0zfMlT}
Locker locations have multiple lockers of various sizes.
\item
\protect\phantomsection\label{NOQJBCy-UEFHYtMg3xFWw}
Only packages within a locker's dimensions are eligible for locker delivery.
\item
\protect\phantomsection\label{h1CFz9yVeEgHUb88u2d5X}
Customers receive a unique code to open the locker when a package is delivered.
\item
\protect\phantomsection\label{jJxQLUtbOXKED8ZhD2gVz}
Packages are held in a locker for a maximum of three days.
\item
\protect\phantomsection\label{d5FzKns1o2dpooHHG0Xbf}
Package pickup must occur within the three-day window and the locker location's operating hours.
\item
\protect\phantomsection\label{SFbsFlIPu1J8MZrlNIMIB}
After three days, uncollected packages are removed, the locker is released, and a refund is initiated.
\item
\protect\phantomsection\label{dvFlDSrFmCTXmO4ql9C6Q}
Each locker can only be assigned to one customer/package at a time.
\item
\protect\phantomsection\label{qlI58ATCQvsiLKg6yU-yp}
The locker's access code is invalidated once a package is collected.
\item
\protect\phantomsection\label{-rs--MK2tJz2BQpN5tzQR}
Customers can return eligible items by selecting a nearby locker location and receiving a new code for an available locker.
\item
\protect\phantomsection\label{K8uEoZ6K9YJTZrJU2ahnZ}
The logistics team picks up returned items using a unique code, after which the customer is notified and refunded.
\end{enumerate}

\subsection{System actors}\label{T4HM0is7bZjPBeXieQAjr}
\\
The following are the primary and secondary actors who interact with the system.

\begin{itemize}
\item
\protect\phantomsection\label{rBQDw4jlarrt6b4haz8QR}
\textbf{Customer:} Selects locker locations, picks up packages using a code, and initiates returns by dropping off packages.
\item
\protect\phantomsection\label{N1H-8bC0jmQI7LkcrYNkR}
\textbf{Delivery person:} Delivers packages to lockers for customer pickup and collects returned items from lockers.
\end{itemize}

\subsection{Important classes and relationships}\label{d2OjqG3pp03E5-pIDSbOc}
\\
Below are the system design's most important classes, descriptions, and key relationships.

\begin{center}
\begin{tabular}{|p{0.18\linewidth}|p{0.42\linewidth}|p{0.35\linewidth}|}
\hline
\textbf{Class} & \textbf{Description} & \textbf{Important Relationships} \\
\hline
Customer & Represents the user who places orders and initiates returns. & Association with Order and Notification. \\
\hline
DeliveryPerson & Represents the logistics staff who deliver packages and collect returns. & Association with LockerPackage and Notification. \\
\hline
Order & Models a transaction initiated by a customer, containing items and a delivery location. & Composition of Item, Association with Notification. \\
\hline
Package & Encapsulates the physical bundle for delivery or return, with a specific size. & Association with Order. \\
\hline
LockerPackage & A specialized package with locker-specific details like an access code and validity period. & Inheritance from Package, Composition of Locker. \\
\hline
Locker & Represents a single physical locker compartment with a specific size and status. & Aggregation by LockerLocation. \\
\hline
LockerLocation & Models a physical site containing multiple lockers, with operating hours and a location. & Composition of Locker, Aggregation by LockerService. \\
\hline
LockerService & The main class that manages all locker locations and orchestrates the core logic. & Aggregation of LockerLocation. \\
\hline
Notification & Manages user communication (e.g., delivery, return, overdue notices). & Association with Customer and DeliveryPerson. \\
\hline
\end{tabular}
\end{center}


\subsection{Design highlights}\label{LGopN2LHpQzQzFtgKfEIj}
\\
This section covers the overall design approach, the core principles applied, and the key design patterns used in the system.

\subsubsection{Design approach}\label{f2E5xpX3VQ38HN--gtqol}
\\
The system was designed using a bottom-up approach. We started by identifying core entities like \texttt{Locker}, \texttt{Package}, and \texttt{Order}. We then built upon these to model more complex operations such as locker assignment, package handling, and secure access, ensuring the design supports concurrency, scalability, and maintainability.

\subsubsection{Design principles}\label{YFXA2SvHGMtphFsPNn5gz}
\\
The design of the Amazon Locker system implicitly incorporates the SOLID principles to ensure it is robust and maintainable.

\begin{itemize}
\item
\protect\phantomsection\label{ubbXRC2GFz78FD7jAU0JG}
\textbf{Single Responsibility principle (SRP):} Each class has a distinct responsibility. For example, the \texttt{Locker} class manages the state of a single locker, while the \texttt{Notification} class is solely responsible for sending communications to users.
\item
\protect\phantomsection\label{YJMTBqU7cQmWzo1oX48VU}
\textbf{Open/Closed principle (OCP):} The system is open to extension but closed for modification. For instance, the \texttt{LockerPackage} class extends the \texttt{Package} class, adding new functionality for locker-specific interactions without altering the base \texttt{Package} class.
\item
\protect\phantomsection\label{uCm0EeOpWkEF2EnRJKSVn}
\textbf{Liskov Substitution principle (LSP):} Subtypes can be substituted for their base types. A \texttt{LockerPackage} object can be used wherever a \texttt{Package} object is expected, ensuring consistent behavior.
\item
\protect\phantomsection\label{IXfF-VDbSrIjgPt9mxmyt}
\textbf{Interface Segregation principle (ISP):} The system's classes and interactions are designed to be specific. For example, the \texttt{Customer} and \texttt{DeliveryPerson} interact with the \texttt{LockerService} through focused methods (\texttt{requestReturn}, \texttt{deliverPackage}) rather than a single, large interface.
\item
\protect\phantomsection\label{cnGhegMINR8hPulxmFjQL}
\textbf{Dependency Inversion principle (DIP):} High-level modules do not depend on low-level modules; both depend on abstractions. The \texttt{LockerService} depends on abstract concepts of lockers and packages, not their concrete implementations, allowing for flexibility.
\end{itemize}

\subsubsection{Design patterns}\label{L0pnwcnpFb7ffabwS2jIt}
\\
The following design patterns were identified as beneficial for structuring the Amazon Locker system.

\begin{itemize}
\item
\protect\phantomsection\label{VtnZdcHwhHd1fk1sq3O4d}
\textbf{Strategy pattern:} This pattern encapsulates different algorithms for tasks like locker assignment (e.g., finding the closest, best-fit, or least-used locker) and OTP generation, allowing these strategies to be swapped easily.
\item
\protect\phantomsection\label{uI41oJxCaoneAVMx0CTf3}
\textbf{Repository pattern:} This pattern separates the data access and business logic. For example, a \texttt{LockerRepository} would handle all database operations for \texttt{Locker} objects, simplifying the \texttt{LockerService}.
\end{itemize}

\subsection{Object interactions}\label{CmW5Tvoj7X-0tz_gKuy88}
\\
As typical interaction flows illustrate, the system's dynamic behavior is driven by collaborations between key objects to accomplish tasks.

\begin{itemize}
\item
\protect\phantomsection\label{nPrta1THn6u-4xm3fz0TV}
\textbf{Package return:} A \texttt{Customer} initiates a return via the \texttt{LockerService}. The service verifies the request, finds an available \texttt{Locker}, and uses the \texttt{Notification} object to send the customer an access code. The customer uses this code to place the package in the \texttt{Locker}, and the \texttt{LockerService} updates the locker's state.
\item
\protect\phantomsection\label{0MOdk7XDTasdi3o8PH0H1}
\textbf{Package pickup:} A \texttt{DeliveryPerson} places a package in a \texttt{Locker}. The \texttt{LockerService} records this, generates a code, and sends a \texttt{Notification} to the \texttt{Customer}. The \texttt{Customer} uses the code at the \texttt{LockerLocation}, which the \texttt{LockerService} validates to open the correct \texttt{Locker}.
\item
\protect\phantomsection\label{VOZyfwSv4PmMO4p2vVrpR}
\textbf{Code validation:} For any locker access, the \texttt{LockerService} receives a code, validates it against the \texttt{LockerPackage} data, and if correct, interacts with the specific \texttt{Locker} object to unlock it.
\end{itemize}

\subsection{System workflow}\label{NlApekhQz-kA9cvPxNvHz}
\\
The following are a few primary end-to-end scenarios that highlight the system's life cycle from a user's perspective:

\begin{itemize}
\item
\protect\phantomsection\label{p_YXYOZvRrxM0FyZPYqae}
\textbf{Order and pickup workflow:} A customer places an order and selects a \texttt{LockerLocation}. A \texttt{DeliveryPerson} delivers the package to an assigned locker. The Amazon\ Locker\ System sends the customer a notification with a pickup code. The customer goes to the locker, enters the code, the system validates it, opens the locker, and the customer retrieves their product.
\item
\protect\phantomsection\label{piPgknuyZtNZ5yo6goqKj}
\textbf{Return workflow:} A customer submits a return request through the system. The system approves it and finds a suitable locker. It then sends a notification with a new code to the customer. The customer goes to the location, enters the code to open the locker, and places the return package inside. A \texttt{DeliveryPerson} is later notified to collect the returned item.
\end{itemize}

% End of section content